\section{Results}\label{sec:results}
All results are REPLAY or FAULT-INJECTED unless stated; no LLM ran. Sessions within a condition are dependent (Sect.~\ref{sec:method}).

\subsection{Coordination and ablations (XA, RQ1, RQ3)}
\begin{table}[tbp]
\centering\scriptsize\setlength{\tabcolsep}{3pt}
\caption{XA, 160 replay sessions per condition (mean per session unless counts). compl./auth./trans.: numbers of the 160 sessions that completed, kept committed scores authoritative, kept valid transitions. chg./rev.: level changes and reversals; agr.: share of RL-phase decisions equal to the PPO (or heuristic) action; gfire: guardrail firings; strat.: decisions whose final source is the Strategy Agent; byp.: sessions with a guardrail bypass; p50/p95: pooled per-turn latency in ms; calls: agent calls.}\label{tab:xa}
\resizebox{\textwidth}{!}{\begin{tabular}{lcccrrrrrrrrr}
\toprule
Cond. & compl. & auth. & trans. & chg. & rev. & agr. & gfire & strat. & byp. & p50 & p95 & calls\\
\midrule
B0 & 160 & 160 & 160 & 15.119 & 11.744 & 0.626 & 9.588 & 0.0 & 0 & 0.546 & 2.056 & -- \\
B7 & 160 & 160 & 160 & 2.675 & 1.256 & -- & 3.831 & 0.0 & 0 & 0.463 & 1.059 & -- \\
B0c & 160 & 160 & 160 & 15.119 & 11.744 & 0.626 & 9.588 & 0.0 & 0 & 0.553 & 2.004 & -- \\
M1 & 160 & 160 & 160 & 12.45 & 9.869 & 0.4 & 6.838 & 6.669 & 0 & 1.589 & 4.411 & 47.68 \\
M2 & 160 & 160 & 160 & 15.119 & 11.744 & 0.626 & 9.588 & 0.0 & 0 & 1.218 & 3.513 & 23.25 \\
M3 & 160 & 160 & 160 & 12.45 & 9.869 & 0.4 & 6.838 & 6.669 & 0 & 1.118 & 3.189 & 23.84 \\
M4 & 160 & 160 & 160 & 2.225 & 1.0 & 0.014 & 0.0 & 19.619 & 3 & 1.573 & 4.357 & 47.68 \\
M5 & 160 & 160 & 160 & 12.45 & 9.869 & 0.4 & 6.838 & 6.669 & 0 & 1.581 & 4.553 & 47.68 \\
M6 & 160 & 160 & 160 & 12.45 & 9.869 & 0.4 & 6.838 & 6.669 & 0 & 1.592 & 4.46 & 47.68 \\
M8 & 160 & 160 & 160 & 12.45 & 9.869 & 0.4 & 6.838 & 6.669 & 0 & 1.575 & 4.35 & 47.68 \\
M7 & 160 & 160 & 160 & 2.706 & 1.312 & -- & 3.812 & 9.037 & 0 & 1.502 & 2.968 & 47.68 \\
M1c & 160 & 160 & 160 & 12.45 & 9.869 & 0.4 & 6.838 & 6.669 & 0 & 1.573 & 4.534 & 47.68 \\
\bottomrule
\end{tabular}
}
\end{table}
\begin{table}[tbp]
\centering\scriptsize\setlength{\tabcolsep}{3pt}
\caption{XA paired differences (A$-$B) over the same 160 sessions: share of sessions with a different difficulty path, paired mean differences with 95\% percentile-bootstrap intervals, and the pooled-turn median latency difference. Intervals are descriptive (dependent sessions).}\label{tab:xapaired}
\resizebox{\textwidth}{!}{\begin{tabular}{lrcccc}
\toprule
A$-$B & path$\ne$ & $\Delta$ level chg. & $\Delta$ reversals & $\Delta$ PPO agree & $\Delta$ p50 ms\\
\midrule
M1$-$B0 & 1.0 & $-$2.669 [$-$3.188,$-$2.131] & $-$1.875 [$-$2.275,$-$1.462] & $-$0.226 [$-$0.244,$-$0.208] & +1.043 [+1.000,+1.079] \\
M2$-$B0 & 0.0 & 0 & 0 & 0 & +0.672 [+0.646,+0.698] \\
M3$-$M1 & 0.0 & 0 & 0 & 0 & $-$0.471 [$-$0.511,$-$0.425] \\
M4$-$M1 & 1.0 & $-$10.225 [$-$10.544,$-$9.900] & $-$8.869 [$-$9.275,$-$8.463] & $-$0.387 [$-$0.403,$-$0.370] & $-$0.016 [$-$0.065,+0.033] \\
M5$-$M1 & 0.0 & 0 & 0 & 0 & $-$0.008 [$-$0.059,+0.037] \\
M6$-$M1 & 0.0 & 0 & 0 & 0 & +0.004 [$-$0.030,+0.053] \\
M8$-$M1 & 0.0 & 0 & 0 & 0 & $-$0.013 [$-$0.057,+0.032] \\
M7$-$B7 & 0.025 & +0.031 [+0.000,+0.081] & +0.056 [+0.006,+0.125] &  & +1.039 [+0.996,+1.070] \\
B0c$-$B0 & 0.0 & 0 & 0 & 0 & +0.008 [$-$0.019,+0.027] \\
M1c$-$M1 & 0.0 & 0 & 0 & 0 & $-$0.016 [$-$0.055,+0.032] \\
\bottomrule
\end{tabular}
}
\end{table}
Every condition completed all 160 sessions, kept every committed score equal to the recorded evaluator output, and kept every level in $[1,5]$ with RL-phase steps of at most one (A1--A3; exact 95\% lower bound for 160/160: 0.977~\cite{cp34}). Removing the Strategy Agent reproduced the original workflow's difficulty, score and question paths in every session (A4), and removing coaching, communication, grounding or authority enforcement reproduced the full configuration's paths (A5): without faults these components do not influence difficulty or probes, as designed.

With both agents the difficulty path differed from the original in all 160 sessions (A7). Level changes fell from 15.12 to 12.45 per session ($-2.67$, interval $[-3.19,-2.13]$) and reversals from 11.74 to 9.87 ($-1.88$, $[-2.28,-1.46]$); agreement with the PPO action fell from 0.626 to 0.400; sessions contained on average 0.59 more turns ($[+0.34,+0.85]$) because proactive probes added follow-ups. In 11.9 decisions per session a Strategy proposal superseded the PPO action before the guardrail layer. The session-level intervals above treat the 160 dependent sessions as independent. A post hoc cluster bootstrap over the 8 recorded-answer sessions gives wider intervals: level changes \phchgmean{} \phchgci{}, reversals \phrevmean{} \phrevci{}, agreement with the PPO action \phagrmean{} \phagrci{} (all 8 clusters share the sign), and extra turns \phfumean{} \phfuci{}, where only \phfusame{} of the 8 recorded sessions shows an increase, so the extra-turn effect is driven by a single recorded session and should not be generalised. In a post hoc leave-one-cluster-out analysis the level-change interval excludes zero when the opposing set is dropped ($-3.76$, $[-4.81,-2.78]$) but includes zero when any one of the other seven is dropped, so the sign is better supported than the interval; the agreement decrease excludes zero in all eight drops. Level changes and reversals decreased in 7 of the 8 answer sets; the exception is the set behind all extra turns. Across the 160 sessions the Strategy Agent emitted \scTotal{} proposals: \scRemEight{} \emph{remediate} (confidence 0.80), \scRecEight{} and \scRecSeven{} \emph{recover} (0.80 and 0.70) and \scSteady{} \emph{steady} (0.45); it never emitted a \emph{stretch} proposal (\scStretch{}), so in these experiments it only held or lowered difficulty and the stretch rule is unexercised. These are stability and coordination measures; fewer reversals are not evidence of better interviewing, and the replay candidate cannot reveal whether the agents' proposals were pedagogically appropriate.

Two results were not anticipated. (i) \emph{Arbitration is not only a precedence rule.} Without it (M4), steady proposals of confidence 0.45 replaced the PPO action almost everywhere: level changes fell to 2.2 per session, agreement with PPO to 0.014, guardrails never fired, and 3 sessions contained a final action the guardrail layer would still override; M4 removes four checks at once, so this does not by itself identify the threshold; XC (below) isolates it. (ii) With PPO unavailable the agents changed the difficulty path in only 4 of 160 sessions (2.5\%; A9 contradicted): on this replay the Strategy Agent mostly agrees with the labelled non-RL heuristic. Canonical guardrails changed no path (A10 contradicted): replay has no audio, so confidence and hesitation stay at their defaults and the divergent regions of the two implementations (G2, G6) are never reached; the 11.8\% grid disagreement says nothing about their frequency in use.

\emph{Cost.} Pooled per-turn latency rose from \xaBzeroMed{} to \xaMoneMed{}\,ms at the median and from \xaBzeroHigh{} to \xaMoneHigh{}\,ms at the 95th percentile (median difference $+1.04$\,ms, $[+1.00,+1.08]$), 47.7 agent calls per session. The agents are in-process rules, so there are no model calls and no token counts; an LLM-backed agent would add its own latency, which we did not measure.

\begin{table}[tbp]
\centering\scriptsize\setlength{\tabcolsep}{3pt}
\caption{Registered predictions and outcomes. XA (above) and XB (below).}\label{tab:preds}
\begin{tabular}{lp{7.6cm}l}
\toprule
ID & Prediction & Outcome\\
\midrule
A1 & completion 100\% in every condition & Supported \\
A2 & committed scores equal the recorded evaluator output in every session of every condition & Supported \\
A3 & difficulty levels stay within [1,5], and every RL-phase decision moves by at most one level, in every session of every condition & Supported \\
A4 & M2 (no Strategy Agent) reproduces B0 (difficulty, score and question paths) in every session & Supported \\
A5 & M3, M5, M6, M8 reproduce M1's difficulty/score/question paths in every session (fault-free) & Supported \\
A6 & M4 (no arbitration) has guardrail\_bypass>0 in >=1 session; M1 has 0 in all & Supported \\
A7 & M1 differs from B0 in difficulty path in >=50\% of sessions & Supported \\
A8 & median per-turn latency M1 - B0 < 50 ms (pooled turns) & Supported \\
A9 & with PPO unavailable, M7 differs from B7 in difficulty path in >=50\% of sessions & Contradicted \\
A10 & canonical guardrails change the difficulty path in >=1 session vs legacy (B0c vs B0) & Contradicted \\
\bottomrule
\end{tabular}
\\[2mm]
\begin{tabular}{lp{7.6cm}l}
\toprule
ID & Prediction & Outcome\\
\midrule
B1 & variant=full: every fault cell completes, keeps committed scores authoritative and transitions valid; no guardrail bypass, no ungrounded leak & Supported \\
B2 & forbidden\_key: full -> authority\_violation>0 and schema\_violation==0; no\_authority -> authority\_violation==0, schema\_violation>0; containment 100\% in both & Supported \\
B3 & adversarial\_valid: no\_arbitration has guardrail\_bypass>0 in >=1 session; full has 0 & Supported \\
B4 & ungrounded coaching: no\_grounding leaks (>=1 session); full has 0 leaks and rejects plans & Supported \\
B5 & strategy timeout (full): pooled p95 per-turn latency exceeds the clean M1 p95 by >=100 ms & Supported \\
B6 & at fault rate 0.3, >=80\% of sessions see the per-agent breaker open at least once (crash/malformed/stale/nan cells) & Contradicted \\
\bottomrule
\end{tabular}

\end{table}

\subsection{Agent faults and safeguards (XB, RQ2)}
\begin{table}[tbp]
\centering\scriptsize\setlength{\tabcolsep}{3pt}
\caption{XB, 40 sessions per fault and variant (160 per fault over four variants). inj.: injected faulty or harmful outputs in the full variant; compl./auth./trans.: sessions completed, with authoritative scores, with valid transitions; bypass: sessions with a guardrail-overridable final action (full / arbitration removed); leak: sessions in which an ungrounded concept was emphasised (full / grounding removed); brk.: sessions in which the per-agent breaker opened (full); p95: pooled per-turn latency.}\label{tab:xb}
\resizebox{\textwidth}{!}{\begin{tabular}{lrcccccrr}
\toprule
Fault & inj. & compl. & auth. & trans. & bypass & leak & brk. & p95\\
 & (full) & \multicolumn{3}{c}{all 4 variants} & full/no-arb. & full/no-gr. & full & ms\\
\midrule
coach:crash & 154 & 160/160 & 160/160 & 160/160 & 0 / 1 & 0 / 0 & 29 & 4.09 \\
coach:malformed & 154 & 160/160 & 160/160 & 160/160 & 0 / 1 & 0 / 0 & 29 & 4.85 \\
coach:stale & 154 & 160/160 & 160/160 & 160/160 & 0 / 1 & 0 / 0 & 29 & 4.33 \\
coach:ungrounded & 950 & 160/160 & 160/160 & 160/160 & 0 / 1 & 0 / 40 & 0 & 4.45 \\
strat:adversarial\_valid & 930 & 160/160 & 160/160 & 160/160 & 0 / 38 & 0 / 0 & 0 & 4.5 \\
strat:crash & 152 & 160/160 & 160/160 & 160/160 & 0 / 34 & 0 / 0 & 29 & 3.83 \\
strat:forbidden\_key & 152 & 160/160 & 160/160 & 160/160 & 0 / 34 & 0 / 0 & 29 & 3.86 \\
strat:malformed & 152 & 160/160 & 160/160 & 160/160 & 0 / 34 & 0 / 0 & 29 & 3.62 \\
strat:nan\_confidence & 152 & 160/160 & 160/160 & 160/160 & 0 / 34 & 0 / 0 & 29 & 4.14 \\
strat:stale & 152 & 160/160 & 160/160 & 160/160 & 0 / 34 & 0 / 0 & 29 & 4.04 \\
strat:timeout & 152 & 160/160 & 160/160 & 160/160 & 0 / 34 & 0 / 0 & 29 & 154.38 \\
\bottomrule
\end{tabular}
}
\end{table}
\textbf{Containment is at ceiling and does not discriminate the safeguards.} In every fault cell and variant all sessions completed with authoritative scores and valid transitions (B1; Table~\ref{tab:xb}); the hub's own bounds and authoritative scoring already enforce these properties, so the safeguards matter for other properties. \textbf{Authority check.} For a forbidden key (152 injections) the full configuration classified all as authority violations and none as schema violations; with enforcement removed all 152 became schema violations and containment was unchanged (B2): the strict schema alone contains such output, and the authority check is classification and defence in depth. \textbf{Arbitration.} A valid proposal opposing the evidence (930 injections) produced final actions the guardrails would override in 38 of 40 sessions without arbitration and in none with it (B3); without arbitration the guardrail layer is skipped entirely, so the same bypass appeared in 34 of 40 sessions of the Strategy fault cells where proposals were dropped or steady. In the full configuration the adversarial proposal was nevertheless the final source of 10.8 decisions per session on average (a guardrail may override it afterwards): a well-formed, confident, evidenced proposal that is wrong cannot be detected by these checks, only bounded by the guardrails and the one-step limit. \textbf{Grounding.} All 950 ungrounded plans were rejected in the full configuration; without the check an ungrounded concept was emphasised in 40 of 40 sessions (B4). \textbf{Timeouts.} A timing-out Strategy Agent raised the pooled p95 per-turn latency from 4.4 to 154.4\,ms, about the configured 150\,ms timeout (B5). \textbf{Breaker.} With a fault rate of 0.3 the per-agent breaker opened in 29 of 40 sessions (72.5\%; the same count in every fault type, because all share one seeded fault schedule per session) and, with a 60\,s cooldown longer than a session, silenced the agent for the rest of that session; the registered threshold of 80\% was not met (B6 contradicted), and the cost of the breaker design is lost agent availability after one burst.

\subsection{Threshold decomposition and sensitivity (XC)}
\begin{table}[tbp]
\centering\scriptsize\setlength{\tabcolsep}{3pt}
\caption{XC (REPLAY, 160 sessions per row, mean per session): only $\tau$ changes; evidence requirement, one-step bound and guardrails stay in force. chg./rev.: level changes and reversals; agr.: share of RL-phase decisions equal to the PPO action; strat.: decisions whose final source is the Strategy Agent; gfire: guardrail firings; byp.: sessions with a guardrail bypass. Cluster-bootstrap intervals over the 8 recorded-answer sessions are in the released summary file.}\label{tab:xc}
\begin{tabular}{rcccccc}\toprule
$\tau$ & chg. & rev. & agr. & strat. & gfire & byp.\\
\midrule
0.0 & 2.225 & 1.000 & 0.014 & 16.075 & 3.544 & 0 \\
0.5 & 12.450 & 9.869 & 0.400 & 6.669 & 6.838 & 0 \\
0.6 & 12.450 & 9.869 & 0.400 & 6.669 & 6.838 & 0 \\
0.7 & 12.450 & 9.869 & 0.400 & 6.669 & 6.838 & 0 \\
0.8 & 12.450 & 9.819 & 0.404 & 6.594 & 6.863 & 0 \\
0.9 & 15.119 & 11.744 & 0.626 & 0.000 & 9.588 & 0 \\
\bottomrule
\end{tabular}

\end{table}
Setting $\tau=0$ alone reproduced M4's collapse (level changes \xctauzerochg{}, PPO agreement \xctauzeroagr{}, as in M4) with no guardrail bypass in any session (C1): ungated \emph{steady} proposals of confidence 0.45 replace the PPO action almost everywhere, and the guardrail layer then fires on a smaller share of decisions. The collapse is therefore caused by the missing confidence threshold and not by the removed guardrails, bound or evidence check; the bypass seen in M4 is a separate consequence of skipping the guardrail layer. The full configuration reproduced XA's M1 paths in every session at $\tau=0.6$ (C2, a determinism check) and $\tau=0.5$ gave identical paths (C3). Behaviour was flat for $0.5\le\tau\le0.8$ and the Strategy Agent was silent at $\tau=0.9$ (identical to the original workflow's agreement of 0.626; C4, C5): the realised confidences are only 0.45, 0.70 and 0.80 (Sect.~\ref{sec:results}, XA), so any $\tau$ in $(0.45,0.70]$ admits exactly the same proposals, $\tau=0.8$ excludes only the 0.70 \emph{recover} proposals and $\tau=0.9$ admits none; the formula separates \emph{steady} from evidenced proposals and hardly distinguishes among evidenced ones. All \xcsupported{} registered XC predictions were supported. This shows that the threshold matters; the plateau is a consequence of the two-valued confidence, not evidence that behaviour is robust to the threshold, and it does not show that 0.6 is better than another value for interviewing.

\subsection{Grounding probe, rebuilt (XD)}
\begin{table}[tbp]
\centering\scriptsize\setlength{\tabcolsep}{3pt}
\caption{XD: constructed feedback texts that disagree with their intended label, over applicable payloads, for the lexical validator, on three subsets of the 64 recorded evaluator payloads.}\label{tab:xd}
\begin{tabular}{lccc}
\toprule
Category & ALL (64) & NO\_OVERLAP (47) & RELABELLED (64)\\
\midrule
C0 faithful & 17/64 & 0/47 & 0/64 \\
C1 verbatim missing as correct & 0/57 & 0/40 & 0/57 \\
C2 verbatim incorrect as correct & 0/10 & 0/4 & 0/10 \\
C3 narrative correctly explained m & 0/57 & 0/40 & 0/57 \\
C4 truncated paraphrase credit & 40/57 & 40/40 & 57/57 \\
C5 score leak & 0/64 & 0/47 & 0/64 \\
C6 generic fluff & 0/64 & 0/47 & 0/64 \\
\bottomrule
\end{tabular}

\end{table}
All 64 stored payloads contain placeholder concept text, and \xdOverlap{} contain the same placeholder string in both the covered and the missing list. In the original accounting 17 of 64 faithful texts were rejected; on the 47 payloads without overlap and on all 64 with numbered placeholders, faithful texts were accepted in 47/47 and 64/64: the 17 rejections were an artefact of contradictory evidence, not validator behaviour. The correction goes the other way for credit claims: truncated paraphrases of unsupported credit were accepted in 40 of 57 applicable cases originally, but in 40/40 on the overlap-free subset and 57/57 with numbered placeholders (95\% interval for 57/57: [0.937,1.000]). The lexical check rejects verbatim claims, narrative restatements of missing concepts, score leaks and generic text in every applicable case, and accepts paraphrased unsupported credit in every applicable case. The evaluation is deterministic and constructed; it says nothing about LLM outputs.
